\hypertarget{appendix-87d7626c8d37}{}\label{referencestart:appendix-87d7626c8d37}
\section{Coordinate handedness and orientation}\label{appendix-87d7626c8d37-coordinate-handedness-and-orientation}

\textbf{Curriculum depth:} Introduced · \textbf{First seen:} Phase6 · \textbf{Latest development:} Phase6

\begin{allmightroledefinition}

\subsection{Definition and mechanism}\label{appendix-87d7626c8d37-definition-and-mechanism}

Handedness is the sign of an ordered basis relative to a chosen orientation. A reflection reverses it; a proper rotation preserves it. It does not specify up axis, camera direction or matrix storage order.

\end{allmightroledefinition}

The sign of \((\mathbf x\times\mathbf y)\cdot\mathbf z\) classifies a nondegenerate frame. Swapping two axes flips the sign.

\subsection{Prerequisites}\label{appendix-87d7626c8d37-prerequisites}

\begin{itemize}
\tightlist
\item
  \hyperref[appendix-2391db8052b1]{Cross products and oriented area} --- Oriented area supplies face normals, frame signs and geometric determinants.
\end{itemize}

\subsection{Used by}\label{appendix-87d7626c8d37-used-by}

Later curriculum sources can deepen this same node; no future coverage is marked complete.

\subsection{Recall and next depth}\label{appendix-87d7626c8d37-recall-and-next-depth}

Trace winding through a mirror while preserving ordinary distances. Full production proofs, target-specific integration and learner mastery require further evidence.

\subsection{Source and navigation}\label{appendix-87d7626c8d37-source-and-navigation}

\href{../02-Mathematics-for-Visual-Computing/main.pdf\#nameddest=01-foundations-i-ch017-oriented-area-tells-us-which-side-of-the-panel-is-front}{Teaching narrative (Mathematics for Visual Computing, Ch. 1)} · \hyperref[appendix-fd090d982a15-phase-6---parts-501-600]{Coverage and source gap} · \hyperref[appendix-a103e7846d21-phase-6---source-analysis-and-verification]{Evidence and limits} · \hyperref[appendix-4ad47bdd269a]{Dependency Map}
